\documentclass[11pt]{article}
\usepackage[margin=1in]{geometry}
\usepackage{amsmath,amssymb,amsthm}
\usepackage{enumitem}
\usepackage[hidelinks]{hyperref}
\usepackage{parskip}

\newcommand{\HH}{\mathbb{H}}
\newcommand{\C}{\mathbb{C}}
\newcommand{\Z}{\mathbb{Z}}
\newcommand{\Gam}{\Gamma}
\newcommand{\Res}{\operatorname{Res}}
\newcommand{\Lstar}{L^{*}}
\newcommand{\taup}{\tau_p}

\title{\textbf{Meromorphic periods, finite-contour $L$-functions, and the\\
positive-weight Rademacher problem: context and novelty assessment}}
\author{Working note (collaborator-facing)}
\date{14 June 2026}

\begin{document}
\maketitle

\begin{abstract}
This note records the full context around the de Rham--Betti (DRB) finite-contour
construction of periods, quasi-periods, and $L$-functions for meromorphic modular
forms, the original motivating goal (a positive-weight Rademacher-type coefficient
formula), the precise analytic obstruction, and an end-to-end novelty assessment
against the locally-harmonic-Maass-form (LHMF) / cycle-integral literature. The
headline conclusions: (i) the ``drag-the-contour-and-sum-residues'' route reconstructs
a known, divergent object for weight $\ge 0$; (ii) the Duncan--Frenkel and DRB-$s$
regularizations are both the wrong regulator for that divergence; (iii) the
wall-crossing/chamber phenomenon is occupied ground (Bringmann--Kane--Kohnen;
L\"obrich--Schwagenscheidt; Choie--Zagier), but the finite-contour extraction of the
completed all-$s$ $L$-function together with de Rham--Betti quasi-periods appears
transverse to that literature and is the defensible novel core.
\end{abstract}

\section{The research program and its goal}

The overarching aim is to relax the standard restriction ``poles only at the cusps''
to genuine interior poles in the upper half-plane $\HH$, and to develop the analytic
number theory (coefficients, $L$-functions, periods) of meromorphic modular forms with
such poles.

The classical template is \textbf{Hardy--Ramanujan 1919} (the multiple-interior-poles
result, distinct from their more frequently cited circle-method work on weakly
holomorphic forms such as the partition generating function). HR~1919 produces exact,
rapidly convergent expansions for the Fourier coefficients of forms like $1/E_6$, which
is meromorphic of weight $-6$ with a pole at $\tau=i$ (where $E_6(i)=0$) and at all
$\Gam$-translates of $i$; the analogous $1/E_4$ has a pole at $\rho=e^{\pi i/3}$. These
are properly meromorphic, interior-pole objects.

The program comprises:
\begin{itemize}[leftmargin=1.4em]
\item \textbf{The 1806 paper} (\texttt{arXiv:1806.09874}): $L$-functions for meromorphic
modular forms regular at the cusps, built from a pole-subtraction (polylogarithm)
step plus a deformed-Mellin contour. It explicitly flags the positive-weight
Rademacher-coefficient problem as open.
\item \textbf{The DRB note} (de Rham--Betti / Brown framework): a finite-contour cocycle
construction yielding the periods $\omega^{\pm}(f)$, the quasi-periods $\eta^{\pm}(f)$,
and an all-$s$ completed $L$-function $\Lstar(f,s)$ for $\Delta$, for the weakly
holomorphic $\widehat\Delta$, and for general meromorphic $f$, with explicit
wall-crossing across interior poles.
\end{itemize}

\section{The objects and the drafts}

\paragraph{The BFK $L$-function (the object DRB reproduces).}
``BFK'' is Bringmann, Fricke, and Kent, \emph{Special $L$-values and periods of weakly
holomorphic modular forms}, Proc.\ Amer.\ Math.\ Soc.\ \textbf{142} (2014), 3425--3439
(the user dates the preprint $\approx$2011). It associates a regularized,
incomplete-gamma completed $L$-series to weakly holomorphic modular forms and reads
critical $L$-values off the error of modularity of an Eichler integral. It is widely
cited ($\sim 50$), but---per the user's own audit of its citation tree---no citing work
performs a \emph{finite} (non-cusp, regularization-free) extraction of that $L$-function.

\paragraph{What DRB does to it.}
DRB integrates $f$ against the polynomial kernel along a \emph{finite} contour pair
between interior basepoints and reads the cocycle off the result. The $S$-segment runs
$-1/\tau_0 \to \tau_0$, the $T$-segment runs $\tau_0-1 \to \tau_0$; at the elliptic
basepoint $\tau_0=i$ the $S$-segment collapses and a single compact integral on
$[i-1,i]$ carries everything. The construction (a) reproduces the BFK completed
$\Lstar(f,s)$ term by term in the Fourier expansion at every $s\in\C$, with numerical
agreement to $\sim 10^{-53}$ for $\Delta$ and $\sim 10^{-47}$ for $\widehat\Delta$;
(b) yields Brown's periods $\omega^{\pm}(\Delta)$ and quasi-periods $\eta^{\pm}(\Delta)$
from a finite endpoint contour with no cusp divergence; and (c) makes ``finiteness
structural'': every integral is between interior points avoiding singularities, so
there is nothing to regulate. The claimed value-add over BFK is therefore a
\emph{proof of correctness plus an explicit, numerically convergent reconstruction},
where BFK had a working-but-unexplained formula.

\section{The original target and the analytic obstruction}

\paragraph{Target.}
A convergent, Rademacher-type exact formula for the Fourier coefficients $c_f(n)$ of a
\emph{positive-weight} meromorphic modular form with interior poles, regular at the
cusp---the natural generalization of HR~1919 past negative weight.

\paragraph{Status of the negative-weight case (solved).}
For weight $<0$ the orbit sum converges and the problem is settled: Petersson's
meromorphic Poincar\'e series; Bringmann--Kane's polar harmonic Maass forms give
Hardy--Ramanujan-type expansions for all negative-weight meromorphic cusp forms bounded
toward $i\infty$; HR~1919 itself handles $1/E_6$.

\paragraph{The crux obstruction (residue scaling).}
The literal ``drag the contour to the cusp and sum the residues over the whole orbit''
move reconstructs a Poincar\'e/Rademacher series whose convergence is governed by the
modular transformation of residues. For $f$ of weight $k$ with a simple pole at $\taup$,
the residue at an image transforms as
\[
  \Res_{\gamma\taup}(f) \;=\; (c\taup+d)^{\,k-2}\,\Res_{\taup}(f),
  \qquad \gamma=\begin{pmatrix} a & b\\ c & d\end{pmatrix}.
\]
Since $\#\{(c,d): |c\taup+d|\le X\}\sim X^2$, the orbit sum $\sum_{(c,d)}|c\taup+d|^{k-2}$
converges \emph{iff} $k<0$ and diverges for every $k\ge 0$. A weight-$4$ example is thus
on the wrong side of the line; HR's $1/E_6$ (weight $-6$) is on the easy side. This is
exactly the obstruction named in 1806 (``when $k\ge1$ they do not converge at all'').

\paragraph{Two convergence problems, not one.}
The polylogarithm subtraction (subtract $\sum_{\taup}\tfrac1{w_Q}\mathrm{Li}_{1-p}(e(\tau-\alpha_Q))$
at each pole image) cures the \emph{lower-strip $q$-series} convergence; it does nothing
about the \emph{orbit-sum} growth $|c\taup+d|^{k-2}$. Conflating the two is the source of
the temptation. The weight-$0$ borderline is exactly where Bringmann--Kane's
``problem of Petersson'' lives: the naive Poincar\'e series diverges and must be
continued by Hecke's trick, ``further complicated by the fact that the Fourier expansion
does not converge everywhere due to singularities in the upper half-plane.''

\section{Two escape routes and their verdicts}

\paragraph{(a) The DRB $s$-variable as a Hecke-trick regulator --- does not work as posed.}
Hecke's trick tames the orbit sum by inserting a damping factor $|c\tau+d|^{-2\nu}$ and
continuing $\nu\to0$. The DRB kernel supplies $\taup^{\,s-1}$ (and the Hurwitz-zeta
$T$-kernel), which controls the \emph{vertical position} of the pole image, not the
\emph{automorphy factor} $|c\taup+d|$ driving the divergence. So $s$ is the wrong
regulator: it does not damp the $(c,d)$ sum. (One genuine subtlety: the residue phases
$(c\taup+d)^{k-2}$ oscillate, so conditional summation / analytic continuation is not a
priori hopeless---but the continuation that works is Hecke's trick, not the $s$-variable.)

\paragraph{(b) Duncan--Frenkel Rademacher-sum regularization --- wrong failure mode.}
Duncan--Frenkel (\emph{Rademacher sums, moonshine and gravity}) regularize sums whose
\emph{seeds are cusp polar parts} ($q^{-m}$ at $i\infty$), with the delicate case being
weight $\le 0$ / conditional convergence. For positive weight their sums are just the
holomorphic Poincar\'e series, absolutely and locally uniformly convergent for weight
$>1$---cusp seed, no interior poles. The divergence here is different in kind: interior-pole
seeds whose summands grow in \emph{modulus} like $|c\taup+d|^{k-2}$ (absolute divergence,
not a conditional/phase problem). Their truncate-and-correct machinery has no purchase on
that growth. The correct precedent is again Hecke's trick \`a la Bringmann--Kane (weight
$0$). Duncan--Frenkel is the right tool for the (already solved) negative-weight side, not
the open positive-weight side.

\section{The inversion, and the LHMF / L\"obrich--Schwagenscheidt comparison}

\paragraph{The inversion.}
Rather than form the divergent orbit sum, take the finite-contour period/$L$-function as
the regularization-free invariant the divergent sum was only ever a proxy for, and let
the wall-crossing residues track the chamber-to-chamber jumps. The question is whether
this is new.

\paragraph{What the literature already contains.}
\emph{Locally harmonic Maass forms} (Bringmann--Kane--Kohnen, IMRN 2015): functions
harmonic except across geodesics, where they jump, becoming continuous only after adding
an explicit \emph{local polynomial} depending on the chamber. \emph{L\"obrich--Schwagenscheidt}
(\texttt{arXiv:2006.10549}, \emph{Locally harmonic Maass forms and periods of meromorphic
modular forms}) construct $\mathcal{H}_{1-k,n}(\tau)$ of weight $2-2k$ as the $n$-th
\emph{period} (in $z$) of Petersson's two-variable kernel
\[
  H_{k,k-1}(z,\tau)=\sum_{M\in\Gam} v^{2k-1}(z-\tau)^{-1}(z-\bar\tau)^{1-2k}\big|_{2-2k,\tau}M,
\]
which in $z$ is a weight-$2k$ meromorphic modular form with poles exactly at the
$\Gam$-translates of $\tau$. Their Theorem~1.3 splits $\mathcal{H}_{1-k,n}$ into
holomorphic and non-holomorphic Eichler integrals plus a locally polynomial part
$\mathcal{P}_{1-k,n}(\tau)$, harmonic off $E_1=\Gam\cdot i\mathbb{R}_+$ (translates of the
imaginary axis), with the chamber-dependent jump given by the \emph{finite} sum
\[
  \frac{i^{1-n}}{2}\!\!\sum_{\substack{M=\left(\begin{smallmatrix}a&b\\c&d\end{smallmatrix}\right)\in\Gam\\ ac>0,\ \operatorname{Re}(M\tau)<0}}\!\!\bigl(\tau^{n}-(-1)^{n}\tau^{2k-2-n}\bigr)\big|_{2-2k}M .
\]
As an application they obtain finite rational formulas for periods of the meromorphic
forms $f_{k,P}$ attached to positive definite binary quadratic forms (poles at CM points,
weight $2k>0$).

\paragraph{Term-by-term comparison with DRB.}
Write the DRB weight as $w$ (so $w=2k$ in LS notation; period-monomial degree $w-2=2k-2$).
\begin{itemize}[leftmargin=1.4em]
\item \textbf{Shared / occupied --- the wall-crossing local polynomial.} DRB's
per-crossing residue jump is $2\pi i\,\sigma\,\Res(f,\taup)\,(X-\taup Y)^{w-2}$. LS's jump
terms are the period-monomials $\tau^n-(-1)^n\tau^{w-2-n}$ slashed by the finitely many
$M$ whose wall has been crossed. These are the \emph{same graded object}: DRB bundles the
monomials $\taup^{\ell}$ into the single generating polynomial $(X-\taup Y)^{w-2}$; LS
exposes them componentwise. The chamber/wall-crossing phenomenon is therefore \emph{not
new to DRB}.
\item \textbf{Weaker-than-hoped: ``general $f$.''} The Petersson kernel $H_{k,k-1}(z,\tau)$
already handles an \emph{arbitrary} single pole orbit (any $\tau$, not just CM), with
higher order accessed by $\tau$-derivatives; multiple orbits are linear combinations. So
generality of the pole configuration is not by itself a strong distinction.
\item \textbf{Transverse --- not present in LS.}
  \begin{enumerate}[label=(\roman*),leftmargin=1.8em]
  \item The completed \emph{all-$s$} $L$-function $\Lstar(f,s)$. LS computes the
  \emph{integer} periods $r_n$ ($0\le n\le 2k-2$)---period-polynomial coefficients---not
  an $L$-function over $s\in\C$.
  \item The \emph{quasi-periods} $\eta^{\pm}$ and the de Rham--Betti period matrix. LS is
  Betti-side periods only; there is no Brown-style quasi-period / de Rham--Betti
  coordinatization.
  \item The \emph{finite contour between interior basepoints}. LS's period is a Mellin
  integral $\int_0^\infty\!\cdots y^n\,dy$ anchored at the cusps (Cauchy principal value on
  walls); DRB never touches the cusp.
  \end{enumerate}
\end{itemize}

\paragraph{Verdict.}
The overlap is essentially confined to the wall-crossing phenomenon, which is
independently known and which the construction itself treats as structural rather than
headline. The finite-contour extraction of the completed $L$-function and the joint
period/quasi-period reconstruction are transverse to the LHMF / cycle-integral line. The
observation that the Bringmann--Schwagenscheidt cycle-integral papers do not cite 1806 is
consistent with this: their deliverable is cycle-integral \emph{periods} of $f_{k,P}$ via
Shintani lifts, a different object from the completed $L$-function of a general
interior-pole meromorphic form. ``Not cited'' is weak evidence on its own; the substantive
argument is the difference in deliverable, and that argument holds.

\section{Novelty ledger}

\paragraph{Occupied (do not claim).}
The chamber / wall-crossing local-polynomial structure for periods of meromorphic forms
(Bringmann--Kane--Kohnen; L\"obrich--Schwagenscheidt; the pole-bearing period
\emph{functions} of Choie--Zagier; the cohomological framework of Bruggeman--Choie--Diamantis).
Negative-weight and weight-$0$ coefficient formulas (Petersson; Bringmann--Kane; HR~1919).
Positive-weight coefficients for the special family $f_{k,D}$ via theta lifts
(L\"obrich--Schwagenscheidt arithmetic; Li--Neururer; Bengoechea).

\paragraph{Plausibly novel (not found pre-empted; must be positioned).}
The finite-contour de Rham--Betti extraction of $\omega^{\pm}$ \emph{and} $\eta^{\pm}$ for
general meromorphic $f$; the all-$s$ $\Lstar(f,s)$ carrying explicit wall-crossing; and the
proof that the BFK incomplete-gamma $L$-function is the correct finite object, with an
explicit numerically convergent reconstruction.

\paragraph{Still open (delivered by none of the above).}
A convergent positive-weight Rademacher \emph{coefficient} formula for $c_f(n)$ of general
meromorphic forms regular at the cusp. The obstruction is the $|c\taup+d|^{k-2}$ growth;
the only known route is Hecke-trick analytic continuation, and it has not been carried out
in this generality.

\section{Priority and prior art (forward-citation sweep)}

A forward-citation sweep of BFK and 1806, filtered for any finite-contour or
de Rham--Betti extraction of the completed $L$-function, returns the following.

\begin{itemize}[leftmargin=1.4em]
\item The completed $L$-function of an interior-pole or weakly holomorphic form is by now
a small but real subfield: BFK (weakly holomorphic); 1806 (meromorphic, regular at the
cusp); L\"obrich--Schwagenscheidt ($L$-values of special meromorphic forms via Cauchy
principal value); Wang--Zhang, \emph{Meromorphic quasi-modular forms and their
$L$-functions} (2022), via regularized integrals; and a completed $L$-function for a
weight-$2$ polar harmonic Maass form (\texttt{arXiv:2201.03146}). Hence ``first
$L$-function of a meromorphic form'' is not an available claim, and $\Lstar(f,s)$ as a
value is the same completed object these define.
\item The 1806 paper \emph{is} cited within this lineage: Wang--Zhang reference it (as
McGady) for ``$L$-functions for meromorphic modular forms holomorphic at infinity.'' By
contrast the LHMF / cycle-integral papers (Bringmann--Kane--Kohnen;
L\"obrich--Schwagenscheidt; Alfes-Neumann--Bringmann--Schwagenscheidt) do not cite it,
consistent with their deliverable being cycle-integral periods rather than the completed
$L$-function. Cited where the object coincides, not cited where it differs---the pattern
is itself evidence of transversality, not oversight.
\item Every paper in the $L$-function niche defines the $L$-function through a
\emph{regularized} Mellin integral from $0$ to $i\infty$ (cusp regularization and/or Cauchy
principal value at poles). None performs the finite-contour cocycle extraction between
interior basepoints; none extracts quasi-periods $\eta^{\pm}$ alongside periods
$\omega^{\pm}$; none frames the result in Brown's de Rham--Betti period structure. The
priority window on the \emph{method} and on the \emph{joint period/quasi-period/$L$-function
package} therefore appears clean.
\end{itemize}

\section{Recommended framing and proposed title}

\textbf{Lead with the $L$-function.} The distinctive content is \emph{uniformity}: a single
construction whose completed $L$-function is the same object across cuspidal, weakly
holomorphic, and interior-pole forms---regimes that previously required mutually incompatible
tools. BFK's regulator is regime-locked (it leans on the harmonic-Maass / non-holomorphic
completion and has no purchase on interior poles); the finite-contour deformation handles
interior poles and also processes $\widehat\Delta$ with the identical kernel; Brown's cocycle
yields $\omega^{\pm},\eta^{\pm}$ but no $L$-function. Nobody has run all of this through one
object. The $L$-function is the right lead object precisely because it is what stays fixed
while the input ranges over the three regimes.

\textbf{The claim to make:} a single finite-contour cocycle yielding the completed all-$s$
$L$-function \emph{together with} the periods $\omega^{\pm}$ and quasi-periods $\eta^{\pm}$,
uniformly across cuspidal, weakly holomorphic, and interior-pole forms, with nothing to
regulate.

\textbf{The phrasing to avoid:} ``first $L$-function for both weakly holomorphic and
meromorphic forms.'' Wang--Zhang's regularized integral arguably already spans
cusp-exponential growth and interior poles within one formalism, so they can contest that
exact priority. What they lack is the rest of the bundle---no periods, no quasi-periods, no
cocycle, regularized rather than finite. State the unification as the cocycle bundle, not as
priority over Wang--Zhang on the bare $L$-function.

\textbf{Proposed manuscript title} (replacing the current working title):
\begin{quote}
\emph{Finite-contour cocycles for (quasi-)periods and $L$-functions of meromorphic modular
forms on $\mathrm{SL}_2(\Z)$.}
\end{quote}
``Finite-contour cocycles'' puts the method---the part with no other occupant---up front;
``(quasi-)periods'' surfaces the $\eta^{\pm}$ / de Rham--Betti differentiator while folding in
$\omega^{\pm}$; ``$L$-functions'' stays in the lead per the uniformity argument above;
``on $\mathrm{SL}_2(\Z)$'' is honest level-$1$ scoping that pre-empts the higher-level reflex.
Read maximally, ``meromorphic'' subsumes the weakly holomorphic and cuspidal cases, so the
span lives in the title; the abstract's opening must then make ``one construction, all three
regimes, $L$ and (quasi-)periods together'' load-bearing from the first line, or a skimmer
files the paper under the occupied $L$-function niche.

\emph{Leaner fallbacks:} \emph{Finite-contour cocycles for (quasi-)periods and $L$-functions
of meromorphic modular forms}; or \emph{(Quasi-)periods and $L$-functions of meromorphic
modular forms via finite-contour cocycles}.

\section{Recommended actions}

\begin{enumerate}[leftmargin=1.6em]
\item Retitle the main manuscript to the proposed title in \S8, and make the uniformity
claim (one cocycle; the $L$-function with $\omega^{\pm},\eta^{\pm}$; three regimes)
load-bearing in the first two sentences of the abstract.
\item Cite Bringmann--Kane--Kohnen (2015) and L\"obrich--Schwagenscheidt (2020) in the DRB
note; they are the closest references and are currently absent.
\item Position the residue jump $(X-\taup Y)^{w-2}$ explicitly against the LS local
polynomial, and test whether, for poles at CM points, the DRB pole-crossing residues are
the specialization of Choie--Zagier rational period functions. This either claims or
disclaims the point-crossing structure cleanly.
\item Frame the contribution as ``finite-contour de Rham--Betti $L$-function and
quasi-period reconstruction'' (transverse to LHMF/LS), \emph{not} as a new wall-crossing
phenomenon and \emph{not} as a positive-weight coefficient formula.
\item If pursuing positive-weight coefficients: the regulator is Hecke's trick (as in the
Bringmann--Kane weight-$0$ Petersson problem), not the DRB $s$-variable and not
Duncan--Frenkel.
\end{enumerate}

\section*{Key references}
\begin{itemize}[leftmargin=1.4em]
\item G.~H.~Hardy, S.~Ramanujan, \emph{Asymptotic formulae in combinatory analysis} and
the 1918/1919 interior-pole work (e.g.\ $1/E_6$).
\item K.~Bringmann, K.-H.~Fricke, Z.~A.~Kent, \emph{Special $L$-values and periods of
weakly holomorphic modular forms}, Proc.\ Amer.\ Math.\ Soc.\ 142 (2014), 3425--3439.
\item [1806] \emph{$L$-functions for meromorphic modular forms} (\texttt{arXiv:1806.09874}).
\item K.~Bringmann, B.~Kane, \emph{Ramanujan and coefficients of meromorphic modular
forms}, J.\ Math.\ Pures Appl.\ (\texttt{arXiv:1603.07079}); \emph{Ramanujan-like formulas
for Fourier coefficients of all meromorphic cusp forms} (\texttt{arXiv:1603.09250}).
\item K.~Bringmann, B.~Kane, \emph{A problem of Petersson about weight $0$ meromorphic
modular forms}, Res.\ Math.\ Sci.\ 3 (2016).
\item K.~Bringmann, B.~Kane, W.~Kohnen, \emph{Locally harmonic Maass forms and the kernel
of the Shintani lift}, IMRN 2015, 3185--3224.
\item S.~L\"obrich, M.~Schwagenscheidt, \emph{Locally harmonic Maass forms and periods of
meromorphic modular forms} (\texttt{arXiv:2006.10549}); \emph{Arithmetic properties of
Fourier coefficients of meromorphic modular forms} (\texttt{arXiv:2010.06297}).
\item W.~Wang, H.~Zhang, \emph{Meromorphic quasi-modular forms and their $L$-functions}
(\texttt{arXiv:2202.09542}, 2022)---completed $L$-functions of meromorphic quasi-modular
forms via regularized integrals; cites 1806 (as McGady).
\item \emph{The completed $L$-function attached to the weight-$2$ polar harmonic Maass form
$H^{*}_{N,z}$} (\texttt{arXiv:2201.03146}).
\item Y.~Choie, D.~Zagier, \emph{Rational period functions for $\mathrm{PSL}_2(\Z)$},
Contemp.\ Math.\ 143 (1993), 89--108; W.~Kohnen, D.~Zagier, \emph{Modular forms with
rational periods} (1984).
\item R.~Bruggeman, Y.~Choie, N.~Diamantis, \emph{Holomorphic automorphic forms and
cohomology}, Mem.\ Amer.\ Math.\ Soc.\ 253 (2018), no.\ 1212.
\item J.~Duncan, I.~Frenkel, \emph{Rademacher sums, moonshine and gravity}, Commun.\
Number Theory Phys.\ 5 (2011), 1--128.
\item F.~Brown, \emph{A class of non-holomorphic modular forms III: real analytic cusp
forms for $SL_2(\Z)$} (\texttt{arXiv:1710.07912}); and personal communication (2019) on
the two-segment cocycle procedure.
\end{itemize}

\end{document}
